\documentclass[aspectratio=169]{beamer}

% =========================================================
% Tema e pacotes
% =========================================================
\usetheme{Madrid}
\usecolortheme{default}

\usepackage[utf8]{inputenc}
\usepackage[T1]{fontenc}
\usepackage[brazil]{babel}
\usepackage{lmodern}
\usepackage{listings}
\usepackage{xcolor}
\usepackage{tikz}
\usepackage{booktabs}
\usepackage{array}
\usepackage{amssymb}
\usepackage{hyperref}
\usetikzlibrary{shapes,arrows.meta,positioning,fit,calc}

% =========================================================
% Cores
% =========================================================
\definecolor{rmiBlue}{RGB}{1,73,124}
\definecolor{rmiTeal}{RGB}{2,128,144}
\definecolor{rmiSea}{RGB}{0,168,150}
\definecolor{rmiLightBlue}{RGB}{225,241,248}
\definecolor{rmiLightGreen}{RGB}{226,245,240}
\definecolor{rmiLightYellow}{RGB}{255,246,210}
\definecolor{rmiGray}{RGB}{245,245,245}
\definecolor{rmiDark}{RGB}{1,73,124}

\setbeamercolor{structure}{fg=rmiBlue}
\setbeamercolor{frametitle}{fg=white,bg=rmiBlue}
\setbeamercolor{title}{fg=white,bg=rmiBlue}
\setbeamercolor{block title}{fg=white,bg=rmiTeal}
\setbeamercolor{block body}{bg=rmiLightBlue}

% =========================================================
% Configuração de codigo Java
% =========================================================
\lstdefinestyle{javastyle}{
    language=Java,
    basicstyle=\ttfamily\scriptsize,
    keywordstyle=\color{blue!75!black}\bfseries,
    commentstyle=\color{green!45!black},
    stringstyle=\color{red!65!black},
    showstringspaces=false,
    breaklines=true,
    frame=leftline,
    framerule=2.5pt,
    rulecolor=\color{rmiTeal},
    backgroundcolor=\color{white},
    xleftmargin=6pt,
    tabsize=4
}
\lstset{style=javastyle}

% =========================================================
% Macros uteis
% =========================================================
\newcommand{\code}[1]{\texttt{#1}}
\newcommand{\ok}{\ensuremath{\checkmark}}
\newcommand{\seta}{\ensuremath{\rightarrow}}

% =========================================================
% Dados da apresentação
% =========================================================
\title{Estudo de Caso: Java RMI}
\subtitle{Remote Method Invocation em Sistemas Distribuídos}
\author{Ana Iara Loayza Costa, José Victor Brito Costa,\\ Hígor Pinheiro Costa, Milena Freire Britto Neves}
\institute{Disciplina de Sistemas Distribuídos}
\date{\today}

\begin{document}

% =========================================================
% SLIDE 1 - CAPA
% =========================================================
\begin{frame}
    \titlepage
\end{frame}

% =========================================================
% SLIDE 2 - ROTEIRO
% =========================================================
% ===============================================
% SLIDE 2: SUMÁRIO
% ===============================================
\begin{frame}{Sumário}
    \tableofcontents
\end{frame}

% =========================================================
% PARTE 1
% =========================================================
\section{Base conceitual}
 
% =========================================================
% SLIDE 3 - PONTE
% =========================================================
\begin{frame}{Da invocação remota ao Java RMI}
    \centering
    \begin{tikzpicture}[
        levelbox/.style={rectangle, rounded corners=6pt, draw=rmiBlue, line width=1.8pt, align=center, minimum width=2.8cm, minimum height=1.2cm, fill=rmiLightBlue},
        arrow/.style={-{Latex[length=4mm, width=3mm]}, line width=2.5pt, rmiTeal}
    ]
        % Sequência horizontal
        \node[levelbox] (inv) at (0, 0) {
            \textbf{Invocação}\\[-2pt]
            \textbf{Remota}\\[4pt]
            \footnotesize Conceito geral
        };
        
        \node[levelbox, fill=white] (rmi) at (4.2, 0) {
            \textbf{RMI}\\[4pt]
            \footnotesize Orientado\\[-2pt]\footnotesize a objetos
        };
        
        \node[levelbox, fill=rmiLightGreen] (java) at (8.4, 0) {
            \textbf{Java RMI}\\[4pt]
            \footnotesize Implementação\\[-2pt]\footnotesize em Java
        };
        
        % Setas de progressão
        \draw[arrow] (inv.east) -- (rmi.west);
        \draw[arrow] (rmi.east) -- (java.west);
        
        % Alternativa: RPC
        \node[levelbox, draw=gray, fill=gray!10] (rpc) at (4.2, -2) {
            \textbf{RPC}\\[4pt]
            \footnotesize Procedural
        };
        \draw[-{Latex[length=4mm, width=3mm]}, line width=1.5pt, gray, dashed] (inv.south) -- ++(0,-0.5) -| (rpc.north);
        
        % Anotações
        \node[font=\footnotesize, align=center, text=rmiDark] at (2.1, 0.9) {especialização};
        \node[font=\footnotesize, align=center, text=rmiDark] at (6.3, 0.9) {implementação};
        \node[font=\footnotesize, align=center, text=gray] at (2.1, -1.35) {alternativa};
    \end{tikzpicture}
 
    \vspace{0.6cm}
    \begin{block}{Ideia Central}
        Java RMI aplica o conceito de invocação remota ao paradigma orientado a objetos.
    \end{block}
\end{frame}

% =========================================================
% SLIDE 4 - O QUE E
% =========================================================
\begin{frame}{O que é Java RMI?}
    \begin{columns}[T]
        \begin{column}{0.50\textwidth}
            \begin{block}{Definicao}
                Java RMI permite que um objeto em uma JVM invoque métodos de um objeto localizado em outra JVM, possívelmente em outra máquina.
            \end{block}

            \vspace{0.2cm}
            \textbf{Em vez de pensar apenas em mensagens, o programador trabalha com:}
            \begin{itemize}
                \item interfaces Java;
                \item objetos remotos;
                \item métodos;
                \item exceções;
                \item serialização.
            \end{itemize}
        \end{column}

        \begin{column}{0.46\textwidth}
            \centering
            \begin{tikzpicture}[
                vm/.style={rectangle, rounded corners, draw=rmiBlue, fill=rmiLightBlue, very thick, align=center, minimum width=3.7cm, minimum height=1.3cm},
                arr/.style={-{Latex[length=3mm]}, very thick, rmiTeal}
            ]
                \node[vm] (a) {JVM A\\\textbf{Cliente}};
                \node[vm, below=1.4cm of a] (b) {JVM B\\\textbf{Objeto remoto}};
                \draw[arr] (a) -- node[right, align=left] {chamada de\\método remoto} (b);

                \draw[arr, dashed, gray] (b.west) -- ++(-0.8,0) |- node[left, pos=0.25, font=\footnotesize, align=right] {retorno} (a.west);
            \end{tikzpicture}

            \vspace{0.5cm}
            \begin{alertblock}{Atenção}
                A sintaxe parece local, mas a semântica é distribuída: existe rede, latência e possibilidade de falha parcial.
            \end{alertblock}
        \end{column}
    \end{columns}
\end{frame}

% =========================================================
% SLIDE 5 - POR QUE
% =========================================================
\begin{frame}[t]{Por que RMI existe?}

    \begin{columns}[T,onlytextwidth]
        % Coluna esquerda
        \begin{column}{0.49\textwidth}
            \textbf{\textcolor{red}{Problema: comunicação em sistemas distribuídos}}

            \begin{itemize}\setlength\itemsep{1pt}
                \item Processos em máquinas diferentes precisam cooperar.
                \item A comunicação direta por sockets é mais verbosa.
                \item O programador precisa lidar com:
                \begin{itemize}\setlength\itemsep{0pt}
                    \item[$\circ$] envio e recebimento de mensagens;
                    \item[$\circ$] formato dos dados;
                    \item[$\circ$] protocolos de comunicação;
                    \item[$\circ$] erros de rede.
                \end{itemize}
            \end{itemize}
        \end{column}

        % Coluna direita
        \begin{column}{0.49\textwidth}
            \textbf{\textcolor{teal}{Solução: proposta do RMI}}

            \begin{itemize}\setlength\itemsep{1pt}
                \item[$\checkmark$] Abstrai parte da comunicação pela rede.
                \item[$\checkmark$] Usa interfaces Java familiares.
                \item[$\checkmark$] Permite invocar métodos remotos.
                \item[$\checkmark$] Mantém o modelo orientado a objetos.
                \item[$\checkmark$] Esconde detalhes de serialização e transporte.
            \end{itemize}
        \end{column}
    \end{columns}

    \vspace{0.3cm}

    \begin{beamercolorbox}[
        rounded=true,
        shadow=false,
        wd=\textwidth,
        sep=0.8ex
    ]{block body}
        \scriptsize
        \textbf{Resumo:}
        RMI não elimina os desafios dos sistemas distribuídos, mas oferece uma
        abstração orientada a objetos para chamar métodos em objetos localizados
        em outra JVM, possivelmente em outra máquina.
    \end{beamercolorbox}
\end{frame}

% =========================================================
% SLIDE 6 - COMPONENTES
% =========================================================
\begin{frame}{Componentes principais do Java RMI}
    \begin{table}
        \small
        \renewcommand{\arraystretch}{1.3}
        \begin{tabular}{>{\bfseries}p{3.2cm}p{9.2cm}}
            \toprule
            Componente & Funcao \\
            \midrule
            Interface remota & Define os métodos que podem ser invocados remotamente. \\
            Objeto remoto & Implementa a interface remota e executa a lógica real no servidor. \\
            Stub ou proxy & Representa o objeto remoto no lado do cliente. \\
            RMI Registry & Serviço de nomes usado para localizar objetos remotos. \\
            Serialização & Empacota parâmetros, retornos e exceções para trafegar pela rede. \\
            RemoteException & Indica falhas especificas de comunicação remota. \\
            \bottomrule
        \end{tabular}
    \end{table}

\end{frame}

% =========================================================
% PARTE 2
% =========================================================
\section{Arquitetura e fluxo}

% =========================================================
% SLIDE 7 - DIVISOR
% =========================================================
\begin{frame}
    \centering
    \Huge{Arquitetura e fluxo de chamada}

    \vspace{0.6cm}
    \large{Como uma chamada de método atravessa a rede?}
\end{frame}

% =========================================================
% SLIDE 8 - ARQUITETURA
% =========================================================
\begin{frame}{Arquitetura geral do Java RMI}
    \centering
    \begin{tikzpicture}[
        node distance=1.1cm,
        block/.style={rectangle, rounded corners, draw=rmiBlue, fill=rmiLightBlue, very thick, align=center, minimum width=2.6cm, minimum height=1.0cm},
        runtime/.style={rectangle, rounded corners, draw=rmiTeal, fill=rmiLightGreen, thick, align=center, minimum width=3.0cm, minimum height=1.0cm},
        registry/.style={rectangle, rounded corners, draw=orange!80!black, fill=rmiLightYellow, thick, align=center, minimum width=3.2cm, minimum height=1.0cm},
        arr/.style={-{Latex[length=3mm]}, very thick, rmiBlue}
    ]
        \node[block] (client) {Aplicacao\\cliente};
        \node[block, right=1.2cm of client] (stub) {Stub\\proxy local};
        \node[runtime, right=1.2cm of stub] (transport) {RMI Runtime\\rede / JRMP};
        \node[block, right=1.2cm of transport] (remote) {Objeto remoto\\servidor};
        \node[registry, below=1.25cm of transport] (reg) {RMI Registry\\nomes \seta{} stubs};

        \draw[arr] (client) -- node[above, font=\scriptsize] {método} (stub);
        \draw[arr] (stub) -- node[above, font=\scriptsize] {marshal} (transport);
        \draw[arr] (transport) -- node[above, font=\scriptsize] {invocação} (remote);
        \draw[arr, dashed] (client.south) .. controls +(0,-1.0) and +(-2.1,0) .. node[below, font=\scriptsize] {lookup inicial} (reg.west);
        \draw[arr, dashed] (remote.south) .. controls +(0,-1.0) and +(2.1,0) .. node[below, font=\scriptsize] {bind/rebind} (reg.east);
    \end{tikzpicture}

    \vspace{0.35cm}
    \begin{block}{Leitura do diagrama}
        O Registry é usado para localizar o objeto remoto. Depois do \code{lookup}, o cliente invoca métodos no stub.
    \end{block}
\end{frame}

% =========================================================
% SLIDE 9 - REGISTRY
% =========================================================
\begin{frame}[fragile]{RMI Registry: serviço de nomes}
    \begin{columns}[T]
        \begin{column}{0.48\textwidth}
            \begin{block}{No servidor}
                \begin{enumerate}
                    \item Cria o objeto remoto.
                    \item Exporta o objeto para receber chamadas.
                    \item Registra o objeto com um nome.
                    \item Aguarda requisicoes.
                \end{enumerate}
            \end{block}

            \begin{lstlisting}[basicstyle=\ttfamily\tiny]
Registry reg = LocateRegistry.createRegistry(1099);
reg.rebind("BancoService", objetoRemoto);
            \end{lstlisting}
        \end{column}

        \begin{column}{0.48\textwidth}
            \begin{block}{No cliente}
                \begin{enumerate}
                    \item Acessa o Registry do servidor.
                    \item Faz \code{lookup} pelo nome.
                    \item Recebé um stub.
                    \item Chama métodos no stub.
                \end{enumerate}
            \end{block}

            \begin{lstlisting}[basicstyle=\ttfamily\tiny]
Registry reg = LocateRegistry.getRegistry("localhost", 1099);
BancoService banco = (BancoService) reg.lookup("BancoService");
            \end{lstlisting}
        \end{column}
    \end{columns}

    \vspace{0.2cm}
    \centering
    \small Porta padrao usual do Registry: \code{1099}.
\end{frame}

% =========================================================
% SLIDE 10 - FLUXO
% =========================================================
\begin{frame}{Fluxo de uma chamada remota}
    \begin{columns}[T]
        \begin{column}{0.52\textwidth}
            \begin{enumerate}
                \item Cliente executa \code{banco.depositar("joao", 250.0)}.
                \item O stub intercepta a chamada.
                \item Parâmetros são empacotados: \textit{marshalling}.
                \item Dados atravessam a rede.
                \item Servidor desempacota os parâmetros.
                \item Objeto remoto executa o método real.
                \item Resultado ou exceção e empacotado.
                \item Resposta retorna ao cliente.
            \end{enumerate}
        \end{column}

        \begin{column}{0.44\textwidth}
            \centering
            \begin{tikzpicture}[
                node distance=0.62cm,
                flowstep/.style={rectangle, rounded corners, draw=rmiBlue, fill=rmiLightBlue, align=center, minimum width=4.6cm, minimum height=0.55cm},
                arr/.style={-{Latex[length=2.4mm]}, thick, rmiBlue}
            ]
                \node[flowstep] (s1) {Cliente chama método};
                \node[flowstep, below=of s1] (s2) {Stub empacota parâmetros};
                \node[flowstep, below=of s2] (s3) {Rede transporta a chamada};
                \node[flowstep, below=of s3] (s4) {Servidor invoca objeto real};
                \node[flowstep, below=of s4] (s5) {Resultado volta ao cliente};
                \draw[arr] (s1) -- (s2);
                \draw[arr] (s2) -- (s3);
                \draw[arr] (s3) -- (s4);
                \draw[arr] (s4) -- (s5);
            \end{tikzpicture}
        \end{column}
    \end{columns}
\end{frame}

% =========================================================
% SLIDE 11 - PARAMETROS
% =========================================================
\begin{frame}[fragile]{Passagem de parâmetros e retorno}
    \begin{columns}[T]
        \begin{column}{0.48\textwidth}
            \begin{block}{Por valor}
                Objetos não remotos são cópiados por serialização.
            \end{block}
            \begin{itemize}
                \item Tipos primitivos e \code{String}.
                \item Objetos que implementam \code{Serializable}.
                \item Alterações na cópia não mudam o original.
            \end{itemize}

\begin{lstlisting}[basicstyle=\ttfamily\tiny]
Pessoa p = servico.buscarPessoa(10);
// Pessoa chega como copia serializada
\end{lstlisting}
        \end{column}

        \begin{column}{0.48\textwidth}
            \begin{block}{Por referência remota}
                Objetos remotos exportados são enviados como stub.
            \end{block}
            \begin{itemize}
                \item O objeto real continua na JVM original.
                \item O receptor recebe uma referência remota.
                \item Permite chamadas de volta ao objeto remoto.
            \end{itemize}

\begin{lstlisting}[basicstyle=\ttfamily\tiny]
servidor.registrarCallback(callback);
// callback remoto e passado como stub
\end{lstlisting}
        \end{column}
    \end{columns}
\end{frame}

% =========================================================
% SLIDE 12 - FALHAS
% =========================================================
\begin{frame}{Chamada remota não e chamada local}
    \begin{columns}[T]
        \begin{column}{0.48\textwidth}
            \textbf{Uma chamada local pode falhar por:}
            \begin{itemize}
                \item erro de programação;
                \item exceção da regra de negocio;
                \item problema de memória;
                \item objeto nulo.
            \end{itemize}
        \end{column}

        \begin{column}{0.48\textwidth}
            \textbf{Uma chamada remota também pode falhar por:}
            \begin{itemize}
                \item servidor fora do ar;
                \item timeout;
                \item perda de conexão;
                \item erro de serialização;
                \item incompatibilidade de classes;
                \item objeto remoto indisponível.
            \end{itemize}
        \end{column}
    \end{columns}

    \vspace{0.35cm}
    \begin{alertblock}{Ponto-chave}
        Por isso, métodos remotos declaram \code{throws RemoteException}. A interface ja avisa que a rede faz parte da chamada.
    \end{alertblock}
\end{frame}

% =========================================================
% PARTE 3
% =========================================================
\section{Implementação}

% =========================================================
% SLIDE 13 - DIVISOR
% =========================================================
\begin{frame}
    \centering
    \Huge{Implementação prática}

    \vspace{0.6cm}
    \large{Exemplo: banco distribuído com interface em tempo real}
\end{frame}

% =========================================================
% SLIDE 14 - PASSOS MODERNOS
% =========================================================
\begin{frame}{Passos para implementar Java RMI}
    \begin{enumerate}
        \item Definir uma \textbf{interface remota} que estende \code{Remote}.
        \item Declarar \code{throws RemoteException} nos métodos remotos.
        \item Criar uma classe que \textbf{implementa} a interface remota.
        \item Exportar o objeto remoto com \code{UnicastRemoteObject}.
        \item Criar ou acessar o \code{RMI Registry}.
        \item Registrar o objeto remoto com \code{bind} ou \code{rebind}.
        \item No cliente, fazer \code{lookup} e invocar métodos no stub recebido.
    \end{enumerate}

    \vspace{0.25cm}
    \begin{alertblock}{Atualização importante}
        Em materiais antigos aparece o passo \code{rmic}. Hoje ele não deve ser tratado como obrigatório: o compilador estático de stubs \code{rmic} foi removido do JDK 15.
    \end{alertblock}
\end{frame}

% =========================================================
% SLIDE 15 - INTERFACE
% =========================================================
\begin{frame}[fragile]{Código 1: interface remota}
\begin{lstlisting}
import java.rmi.Remote;
import java.rmi.RemoteException;
import java.util.List;

public interface BancoService extends Remote {
    void depositar(String conta, double valor)
            throws RemoteException;
    void sacar(String conta, double valor)
            throws RemoteException, SaldoInsuficienteException;
    double consultarSaldo(String conta)
            throws RemoteException;
    List<Transacao> getExtrato(String conta)
            throws RemoteException;
    void registrarListener(String conta, ClienteListener l)
            throws RemoteException;
}
\end{lstlisting}

\vspace{0.2cm}
\begin{itemize}
    \item \code{Remote} marca a interface como remota.
    \item \code{getExtrato} retorna uma lista de objetos — demonstra \textbf{passagem por valor}.
    \item \code{registrarListener} recebe um objeto remoto — demonstra \textbf{passagem por referência remota}.
\end{itemize}
\end{frame}

% =========================================================
% SLIDE 16 - TRANSACAO E EXCECAO
% =========================================================
\begin{frame}[fragile]{Código 2: objetos que trafegam pela rede}
    \begin{columns}[T]
        \begin{column}{0.48\textwidth}
            \begin{block}{Passagem por valor}
\begin{lstlisting}[basicstyle=\ttfamily\tiny]
public class Transacao
        implements Serializable {
    private static final long
        serialVersionUID = 1L;
    private final Tipo tipo;
    private final double valor;
    private final double saldoApos;
    private final LocalDateTime dataHora;
    // getters + toString...
}
\end{lstlisting}
            \end{block}
            \small Chega ao cliente como \textbf{cópia serializada}. Alterações não afetam o servidor.
        \end{column}

        \begin{column}{0.48\textwidth}
            \begin{block}{Exceção serializada}
\begin{lstlisting}[basicstyle=\ttfamily\tiny]
public class SaldoInsuficienteException
        extends Exception
        implements Serializable {
    private final double saldoAtual;
    private final double valorSolicitado;
    // construtor com mensagem formatada
}
\end{lstlisting}
            \end{block}
            \small Criada no servidor, \textbf{serializada, enviada pela rede} e capturada no cliente com \code{catch} normal.
        \end{column}
    \end{columns}
\end{frame}

% =========================================================
% SLIDE 17 - IMPLEMENTACAO DO OBJETO REMOTO
% =========================================================
\begin{frame}[fragile]{Código 3: implementação do objeto remoto}
\begin{lstlisting}
public class BancoServiceImpl
        extends UnicastRemoteObject
        implements BancoService {

    private final Map<String, Double> saldos = new HashMap<>();
    private final Map<String, List<Transacao>> extratos = new HashMap<>();
    private final Map<String, List<ClienteListener>> listeners = new HashMap<>();

    public synchronized void sacar(String conta, double valor)
            throws RemoteException, SaldoInsuficienteException {
        double saldoAtual = saldos.get(conta);
        if (valor > saldoAtual)
            throw new SaldoInsuficienteException(saldoAtual, valor);
        // atualiza saldo, registra transacao, notifica listeners...
    }
}
\end{lstlisting}

\vspace{0.1cm}
\small \code{synchronized} garante que operações concorrentes de múltiplos clientes sejam seguras.
\end{frame}

% =========================================================
% SLIDE 18 - CALLBACK
% =========================================================
\begin{frame}[fragile]{Código 4: callback — passagem por referência remota}
    \begin{columns}[T]
        \begin{column}{0.48\textwidth}
            \begin{block}{Interface de callback}
\begin{lstlisting}[basicstyle=\ttfamily\tiny]
public interface ClienteListener
        extends Remote {
    void onTransacao(Transacao t)
        throws RemoteException;
}
\end{lstlisting}
            \end{block}
            \begin{block}{Cliente exporta o objeto}
\begin{lstlisting}[basicstyle=\ttfamily\tiny]
ClienteListenerImpl cb =
    new ClienteListenerImpl("joao");
banco.registrarListener("joao", cb);
// cb e enviado como stub ao servidor
\end{lstlisting}
            \end{block}
        \end{column}

        \begin{column}{0.48\textwidth}
            \begin{block}{Servidor chama de volta}
\begin{lstlisting}[basicstyle=\ttfamily\tiny]
// dentro de depositar() no servidor:
for (ClienteListener l : listeners) {
    l.onTransacao(t); // chama o cliente!
}
\end{lstlisting}
            \end{block}

            \vspace{0.2cm}
            \begin{alertblock}{Inversão de papéis}
                O \textbf{servidor chama o cliente}. O objeto remoto está na JVM do cliente; o servidor possui apenas um stub dele.
            \end{alertblock}
        \end{column}
    \end{columns}
\end{frame}

% =========================================================
% SLIDE 19 - SERVIDOR E REGISTRY
% =========================================================
\begin{frame}[fragile]{Código 5: servidor e registry}
\begin{lstlisting}
public class Servidor {
    public static void main(String[] args) throws Exception {
        BancoServiceImpl banco = new BancoServiceImpl();

        Registry registry = LocateRegistry.createRegistry(1099);
        registry.rebind("BancoService", banco);

        System.out.println("Servidor RMI pronto na porta 1099.");
    }
}
\end{lstlisting}

\vspace{0.2cm}
\begin{lstlisting}
// No cliente:
Registry registry = LocateRegistry.getRegistry("localhost", 1099);
BancoService banco =
    (BancoService) registry.lookup("BancoService");
// banco e um stub -- parece local, e remoto
\end{lstlisting}

\vspace{0.15cm}
\small O cliente não conhece \code{BancoServiceImpl}. Ele conhece apenas a interface e o stub retornado pelo Registry.
\end{frame}

% =========================================================
% SLIDE 20 - DEMO
% =========================================================
\begin{frame}[fragile]{Demonstração prática}
    \begin{columns}[T]
        \begin{column}{0.48\textwidth}
            \textbf{O que será demonstrado}
            \begin{enumerate}
                \item Servidor Java sobe: RMI \code{:1099} + REST \code{:8080}.
                \item Interface React abre em \code{localhost:5173}.
                \item Selecionar uma conta — saldo exibido em tempo real.
                \item Realizar depósito — feed RMI pisca instantaneamente.
                \item Realizar saque — extrato atualiza automaticamente.
                \item Forçar saldo insuficiente — exceção serializada exibida.
            \end{enumerate}
        \end{column}

        \begin{column}{0.48\textwidth}
            \textbf{Como subir}
\begin{lstlisting}[language=bash,basicstyle=\ttfamily\tiny]
./run.sh
# compila Java, sobe servidor
# e React automaticamente
\end{lstlisting}

            \vspace{0.3cm}
            \begin{alertblock}{O que observar}
                O feed \textit{RMI em tempo real} mostra o callback chegando ao browser via SSE — o servidor chamou o listener no cliente, que propagou para a interface.
            \end{alertblock}
        \end{column}
    \end{columns}
\end{frame}

% =========================================================
% PARTE 4
% =========================================================
\section{Análise crítica}

% =========================================================
% SLIDE 20 - CARACTERISTICAS
% =========================================================
\begin{frame}{Caracteristicas Importantes}
    \begin{columns}[T]
        \begin{column}{0.48\textwidth}
            \begin{block}{Serialização automatica}
                \begin{itemize}
                    \item Parâmetros são empacotados pelo RMI.
                    \item Retornos e exceções também trafegam pela rede.
                    \item Objetos por valor precisam ser serializaveis.
                \end{itemize}
            \end{block}
        \end{column}

        \begin{column}{0.48\textwidth}
            \begin{block}{Garbage Collection Distribuido}
                \begin{itemize}
                    \item O RMI acompanha referências remotas.
                    \item Usa uma ideia de \textit{lease} para referências ativas.
                    \item Ajuda a liberar objetos remotos não referênciados.
                \end{itemize}
            \end{block}
        \end{column}
    \end{columns}

    \vspace{0.3cm}
    \begin{block}{Por que isso importa?}
        Esses recursos mostram que RMI não é apenas uma chamada pela rede: ele integra comunicação remota ao modelo de objetos da plataforma Java.
    \end{block}
\end{frame}

% =========================================================
% SLIDE 21 - SEGURANCA
% =========================================================
\begin{frame}{Segurança e cuidados atuais}
    \begin{itemize}
        \item RMI expoe métodos pela rede; portanto, e necessário controlar quem pode chama-los.
        \item Em ambientes reais, podem ser necessárias autenticação, autorização e canal seguro.
        \item Java oferece fabricas de socket RMI com SSL/TLS, como \code{SslRMIClientSocketFactory} e \code{SslRMIServerSocketFactory}.
        \item Serialização merece cuidado: dados e classes recebidas não devem ser tratados como confiáveis automaticamente.
        \item Materiais antigos citam \code{SecurityManager}; em Java moderno, ele deixou de ser o caminho recomendado e foi desabilitado permanentemente no Java 24.
    \end{itemize}

    \vspace{0.25cm}
    \begin{alertblock}{Mensagem para a apresentação}
        RMI deve ser usado preferêncialmente em redes controladas ou sistemas internos, não como API pública aberta na Internet.
    \end{alertblock}
\end{frame}

% =========================================================
% SLIDE 22 - LIMITACOES
% =========================================================
\begin{frame}{Limitações e desafios}
    \begin{table}
        \small
        \renewcommand{\arraystretch}{1.25}
        \begin{tabular}{>{\bfseries}p{3.4cm}p{4.5cm}p{4.5cm}}
            \toprule
            Desafio & Problema & Cuidado pratico \\
            \midrule
            Latencia & Chamadas remotas são mais lentas que chamadas locais. & Evitar chamadas muito finas; agrupar operações. \\
            Falhas parciais & Cliente, servidor ou rede podem falhar separadamente. & Tratar \code{RemoteException}, timeouts e recuperacao. \\
            Acoplamento & Cliente e servidor compartilham interfaces Java. & Manter contratos estaveis e versionados. \\
            Serialização & Mudancas em classes podem quebrar compatibilidade. & Usar \code{serialVersionUID} e evolucao cuidadosa. \\
            Firewall/NAT & Objetos exportados sem porta fixa podem dificultar configuração. & Definir portas e politicas de rede. \\
            \bottomrule
        \end{tabular}
    \end{table}
\end{frame}

% =========================================================
% SLIDE 23 - COMPARACAO
% =========================================================
\begin{frame}{RMI vs outras tecnologias}
    \begin{table}
        \scriptsize
        \renewcommand{\arraystretch}{1.25}
        \begin{tabular}{>{\bfseries}p{2.2cm}p{4.7cm}p{5.3cm}}
            \toprule
            Tecnologia & Melhor quando... & Observação \\
            \midrule
            Java RMI & Ambiente 100\% Java, rede controlada, sistemas internos ou legados. & Modelo orientado a objetos e forte integracao com Java. \\
            REST/HTTP & APIs web, integracao entre linguagens e serviços publicos. & Mais comum em sistemas modernos e interoperáveis. \\
            gRPC & Alto desempenho, contratos bem definidos e multiplas linguagens. & RPC moderno com Protocol Buffers. \\
            CORBA & Sistemas legados heterogêneos. & Historicamente importante, mas complexo e menos usado hoje. \\
            \bottomrule
        \end{tabular}
    \end{table}

    \vspace{0.2cm}
    \begin{block}{Leitura crítica}
        RMI não e necessariamente ``melhor'' ou ``pior''. Ele faz sentido quando o contexto e Java, interno e orientado a objetos.
    \end{block}
\end{frame}

% =========================================================
% SLIDE 24 - ONDE FAZ SENTIDO
% =========================================================
\begin{frame}{Onde RMI faz sentido hoje?}
    \begin{columns}[T]
        \begin{column}{0.48\textwidth}
            \textbf{Cenários adequados}
            \begin{itemize}
                \item Sistemas internos escritos em Java.
                \item Aplicacoes legadas que já usam RMI.
                \item Prototipacao e ensino de sistemas distribuídos.
                \item Comunicacao entre componentes Java em rede controlada.
            \end{itemize}
        \end{column}

        \begin{column}{0.48\textwidth}
            \textbf{Exemplo real no ecossistema Java}
            \begin{itemize}
                \item JMX remoto pode usar conectores baseados em RMI.
                \item Isso permite monitorar e gerenciar aplicações Java remotamente.
            \end{itemize}

            \vspace{0.35cm}
            \begin{alertblock}{Cuidado}
                Para APIs públicas ou ambientes heterogêneos, REST ou gRPC costumam ser escolhas mais naturais.
            \end{alertblock}
        \end{column}
    \end{columns}
\end{frame}

% =========================================================
% SLIDE 25 - FECHAMENTO
% =========================================================
\begin{frame}{Conclusão}
    \begin{block}{Ideia principal}
        Java RMI é um estudo de caso classico de invocação remota orientada a objetos.
    \end{block}

    \vspace{0.2cm}
    \textbf{O que aprendemos:}
    \begin{itemize}
        \item RMI aplica a ideia de invocação remota ao modelo de objetos Java.
        \item O cliente chama métodos por meio de uma interface remota.
        \item O stub representa o objeto remoto no lado do cliente.
        \item O Registry ajuda a localizar objetos remotos.
        \item Serialização permite transportar parâmetros e resultados.
        \item \code{RemoteException} lembra que chamada remota não e chamada local.
    \end{itemize}

    \vspace{0.2cm}
    \begin{alertblock}{Mensagem final}
        RMI facilita a programação distribuída em Java, mas não elimina os desafios fundamentais: latencia, falhas, acoplamento, versionamento e segurança.
    \end{alertblock}
\end{frame}

% =========================================================
% SLIDE 26 - REFERENCIAS
% =========================================================
\begin{frame}{Referências}
    \scriptsize
    \begin{itemize}
        \item TANENBAUM, A. S.; VAN STEEN, M. \textit{Sistemas Distribuídos: Princípios e Paradigmas}. 2. ed. Pearson, 2007.
        \item COULOURIS, G.; DOLLIMORE, J.; KINDBERG, T.; BLAIR, G. \textit{Sistemas Distribuídos: Conceitos e Projeto}. 5. ed. Bookman, 2013.
        \item ORACLE. \textit{Java Remote Method Invocation Specification}. Disponível em: \url{https://docs.oracle.com/en/java/javase/25/docs/specs/rmi/index.html}
        \item ORACLE. \textit{Java RMI - Distributed Object Model}. Disponível em: \url{https://docs.oracle.com/en/java/javase/17/docs/specs/rmi/objmodel.html}
        \item ORACLE. \textit{Removed Tools and Components - Removal of RMI Static Stub Compiler}. Disponível em: \url{https://docs.oracle.com/en/java/javase/26/migrate/removed-tools-components.html}
        \item OPENJDK. \textit{JEP 486: Permanently Disable the Security Manager}. Disponível em: \url{https://openjdk.org/jeps/486}
        \item ORACLE. \textit{RMIConnectorServer - Java Management RMI}. Disponível em: \url{https://docs.oracle.com/en/java/javase/25/docs/api/java.management.rmi/javax/management/remote/rmi/RMIConnectorServer.html}
        \item WALDO, J. et al. \textit{A Note on Distributed Computing}. Sun Microsystems Laboratories, 1994.
    \end{itemize}
\end{frame}

% =========================================================
% SLIDE 27 - PERGUNTAS
% =========================================================
\begin{frame}
    \centering
    \vspace{1cm}
    {\Huge Perguntas?}

    \vspace{1cm}
    {\Large Obrigado pela atenção!}
\end{frame}

\end{document}
